Au cours de cette première année, nous avons tout d'abord étudié les évènements
to\-po\-lo\-gi\-ques ayant lieu lors de l'application de règles. %
Ces é\-vé\-ne\-ments sont la création d'une entité topologique, sa suppression, son
non-changement, sa modification, la scission de deux entités topologiques et
leur fusion. %

Cette étude a d'abord donné lieu à l'implémentation d'une méthode pour mettre en
évidence les évènements ayant lieu lors de la construction d'un modèle 3D. %
Pour cela nous avons donc analysé des exemples puis documenté et formalisé ce
qui dans leur écriture provoque un évènement pour chaque entité impactée par la
règle. %
Ainsi, nous avons défini de manière générique et pour chaque entité topologique
les conditions dans lesquelles une règle provoquera un évènement donné. %


Après cela, nous avons pu travailler sur des structures de données, que l'on
appelle journaux d'historiques. %
Ceux-ci permettent, à partir d'un nom persistant, de tracer l'origine et
l'évolution de l'entité désignée dans la spécification paramétrique. %
L'étape suivante a consisté à enregistrer les règles utilisées et leurs
paramètres dans une spécification paramétrique puis de développer des structures
de données pour représenter ces journaux d'historiques. %
Nous avons implémenté un algorithme permettant de construire ces journaux
d'historiques. %

Par la suite, nous avons défini une structure de données pour représenter un
arbre d'appariement. %
Construit à partir d'un journal d'historique et d'une spécification
pa\-ra\-mé\-tri\-que, ce dernier permet de retrouver dans le modèle réévalué,
les entités correspondant aux entités désignées dans la spécification
paramétrique initiale. %
En l'état actuel, nous avons ensuite pu implémenter une méthode pour la
reconstruction automatique d'un modèle à deux cond\-i\-t\-i\-ons : il doit
s'agir d'une construction dont la spécification pa\-ra\-mé\-tri\-que n'a pas été
éditée et ses règles ne doivent pas provoquer de fusion. %

À ce jour, nous avons donc pu implémenter l'enregistrement de spécifications
para\-mé\-tri\-ques, la construction de journaux d'historique, d'arbres
d'appariement et le rejeu de modèles non-édités ne présentant pas d'évènement de
fusion. %
Nous avons également un document de travail dans lequel sont formalisés nos
travaux, document servira pour la rédaction du manuscript de thèse et pour la
rédaction d'un futur article. %


D'ici décembre, nous prévoyons d'intégrer les cas de fusion, d'implémenter
l'édition des spécifications paramétriques avec l'ajout, la suppression et le
déplacement d'appels de règles. %
Puis nous enchaînerons en fin d'année sur l'implémentation du rejeu de
spé\-ci\-fi\-ca\-t\-i\-ons paramétriques après édition. %
Les arbres d'appariement permettront ainsi de tracer les évènements et de
proposer des options dans les cas où plusieurs paramètres topologiques
apparaissent ou si, au contraire, aucun paramètre topologique ne peut être
trouvé. %
L'in\-té\-gra\-t\-i\-on d'une dimension temporelle pour animer l'évolution d'un
modèle devra cer\-tai\-ne\-ment être mise de côté par soucis de temps afin d'une
part d'étendre notre travail à l'utilisation d'opérations plus complexes
définies à partir de scripts de règles, et d'autre part d'intégrer au modèle des
annotations historiques. %

%%% Local Variables:
%%% mode: latex
%%% TeX-master: "../main"
%%% End:
